\section{Mathematical Foundations}

\subsection{Problem Fomulation}
% TODO: write content

\subsection{Scale-Invariant Feature Detection and Matching}
To establish correspondences between the UAV imagery and the reference map, we utilize the Scale-Invariant Feature Transform (SIFT) algorithm, which represents keypoints as 128-dimensional continuous feature vectors. The similarity between a query descriptor $v$ and a train descriptor $u$ is determined using the Euclidean distance ($L_2$ norm) in the feature space:
$$d(v, u) = \sqrt{\sum_{k=1}^{128} (v_k - u_k)^2}$$
Due to the highly repetitive nature of aerial landscapes, direct nearest-neighbor matching is prone to false positives. To mathematically mitigate this, correspondences are filtered using Lowe's ratio test. For any given descriptor, let $d_1$ be the distance to its closest neighbor and $d_2$ be the distance to the second closest neighbor. A match is considered reliable and unique only if it satisfies the following inequality:
$$\frac{d_1}{d_2} < \tau$$
where $\tau$ is a predefined ratio threshold.

\subsection{Robust Estimation via RANSAC}

Given a set of noisy correspondences between UAV image points and map points,
\[
\mathcal{D} = \{(x_i, y_i) \leftrightarrow (X_i, Y_i)\}_{i=1}^{N},
\]
the goal is to estimate transformation parameters $\theta$ (affine model)
that are consistent with the largest subset of correct matches (inliers),
while rejecting incorrect correspondences (outliers).

\textbf{Reprojection Error Definition.}

Each correspondence induces a reprojection error under model $\theta$:
\[
d_i(\theta) = \left\| 
\begin{bmatrix} X_i \\ Y_i \end{bmatrix}
-
\hat{T}_\theta
\begin{bmatrix} x_i \\ y_i \\ 1 \end{bmatrix}
\right\|_2,
\]
where $\hat{T}_\theta$ is the affine transformation matrix.

A point is considered an inlier if:
\[
d_i(\theta) < \epsilon,
\]
where $\epsilon$ is a predefined threshold.

Thus, the objective is:
\[
\theta^* = \arg\max_{\theta} \sum_{i=1}^{N} \mathbf{1}\big(d_i(\theta) < \epsilon\big),
\]
i.e., maximize the number of inliers.

\textbf{RANSAC Iterative Procedure.}

At each iteration, a minimal subset of $s$ correspondences is sampled:
\[
\mathcal{S} \subset \mathcal{D}, \quad |\mathcal{S}| = s,
\]
where for the affine model $s = 3$.

Using this subset, a candidate model $\theta$ is estimated by solving
the linear system defined by the affine transformation.

The model is then evaluated on the full dataset by computing
the reprojection errors $\{d_i(\theta)\}_{i=1}^{N}$ and constructing
the inlier set:
\[
\mathcal{I}(\theta) = \{ i \mid d_i(\theta) < \epsilon \}.
\]

The model with the largest consensus set is selected:
\[
\theta^* = \arg\max_{\theta} |\mathcal{I}(\theta)|.
\]

\textbf{Adaptive Number of Iterations.}

The number of RANSAC iterations is derived probabilistically.
Let $w$ denote the probability that a randomly selected correspondence is an inlier:
\[
w = \frac{|\mathcal{I}|}{N}.
\]

The probability that all $s$ points in a randomly sampled subset are inliers is:
\[
P_{\text{all-inliers}} = w^s.
\]

Thus, the probability that a sampled subset contains at least one outlier is:
\[
1 - w^s.
\]

If we perform $k$ independent iterations, the probability that all sampled subsets
contain at least one outlier is:
\[
(1 - w^s)^k.
\]

Therefore, the probability that at least one of the $k$ samples consists entirely
of inliers is:
\[
1 - (1 - w^s)^k.
\]

To guarantee a desired confidence level $p$, we require:
\[
1 - (1 - w^s)^k \geq p.
\]

Rearranging the inequality:
\[
(1 - w^s)^k \leq 1 - p.
\]

Taking logarithms on both sides:
\[
k \log(1 - w^s) \leq \log(1 - p).
\]

Finally, solving for $k$:
\[
k = \frac{\log(1 - p)}{\log(1 - w^s)}.
\]

\textbf{Final Model Estimation.}

After identifying the best inlier set $\mathcal{I}^*$, the final model
parameters are recomputed using all inliers by solving the
overdetermined system:
\[
A v = b,
\]
using the least-squares solution:
\[
v = (A^T A)^{-1} A^T b.
\]



\subsection{Geometric Transformation Models}
Geometric transformations define the mathematical mapping between UAV imagery and the reference map. We consider three linear models:

\begin{itemize}
	\item \textbf{Similarity Transformation:} A 4-degree-of-freedom model accounting for translation, rotation, and uniform scaling.
	\item \textbf{Affine Transformation:} A 6-degree-of-freedom model extending similarity by adding shear and non-uniform scaling.
	\item \textbf{Projective:} An 8-degree-of-freedom model representing the true 3D to 2D projection.
\end{itemize}

The Similarity model requires perfectly level flight and uniform zooming, failing with axial stretching or tilt. The Affine model approximates slight perspective shifts via shear, but only the Projective model handles significant pitch and roll by modeling converging parallel lines.

We chose the Affine model because it flexibly handles non-ideal conditions (slight tilt or shearing) while remaining mathematically simpler and more stable than the Projective model.

For the selected Affine model, the mapping between UAV coordinates $(x, y)$ and map coordinates $(X, Y)$ is:
\begin{equation}
\begin{bmatrix} X \\ Y \\ 1 \end{bmatrix} =
\begin{bmatrix} a_1 & a_2 & t_x \\ a_3 & a_4 & t_y \\ 0 & 0 & 1 \end{bmatrix}
\begin{bmatrix} x \\ y \\ 1 \end{bmatrix}
\end{equation}

During robust fitting, candidate transformations are scored by reprojection error for each correspondence:
\begin{equation}
d = \sqrt{(X_i - \hat{X}_i)^2 + (Y_i - \hat{Y}_i)^2} < \epsilon
\end{equation}

Given the final set of $n$ inliers, parameters are re-estimated from the overdetermined linear system $A\mathbf{v} = \mathbf{b}$ with $\mathbf{v} = [a_1, a_2, t_x, a_3, a_4, t_y]^T$:
\begin{equation}
\begin{bmatrix}
x_1 & y_1 & 1 & 0 & 0 & 0 \\
0 & 0 & 0 & x_1 & y_1 & 1 \\
\vdots & \vdots & \vdots & \vdots & \vdots & \vdots \\
x_n & y_n & 1 & 0 & 0 & 0 \\
0 & 0 & 0 & x_n & y_n & 1
\end{bmatrix}
\begin{bmatrix} a_1 \\ a_2 \\ t_x \\ a_3 \\ a_4 \\ t_y \end{bmatrix} =
\begin{bmatrix} X_1 \\ Y_1 \\ \vdots \\ X_n \\ Y_n \end{bmatrix}
\end{equation}

\subsection{Least-Squares Parameter Estimation}
The optimal solution $\mathbf{v}_{final}$ is the one that minimizes the sum of squared residuals, defined by the cost function $J(\mathbf{v})$:
\begin{equation}
J(\mathbf{v}) = \|A\mathbf{v} - \mathbf{b}\|^2 = (A\mathbf{v} - \mathbf{b})^T(A\mathbf{v} - \mathbf{b})
\end{equation}

Expanding the terms, we obtain:
\begin{equation}
J(\mathbf{v}) = \mathbf{v}^T A^T A \mathbf{v} - 2\mathbf{b}^T A \mathbf{v} + \mathbf{b}^T \mathbf{b}
\end{equation}

To find the minimum, we compute the derivative with respect to $\mathbf{v}$ and set it to zero:
\begin{equation}
\frac{\partial J(\mathbf{v})}{\partial \mathbf{v}} = 2A^T A \mathbf{v} - 2A^T \mathbf{b} = 0
\end{equation}

Solving this linear equation yields the final parameter vector:
\begin{equation}
A^T A \mathbf{v} = A^T \mathbf{b} \implies \mathbf{v}_{final} = (A^T A)^{-1} A^T \mathbf{b}
\end{equation}
